\section{The Masculine and Feminine Races}

\begin{quotationx}

In this chapter of \textit{Sintesi}, \textbf{Julius Evola} demonstrates that the male and female are equivalent to different races. This is not simply the obvious corporeal differences, but extends also the the soul and the spirit. The female spirit is undifferentiated whereas the male spirit is differentiated.

He appeals to \textbf{J J Bachofen}`s writings on ancient gynocentric civilizations which had ``promiscuity, communism, mother right, general equality as their characteristics''. He mentions two possible outcomes:

\begin{itemize}


\item an awakening, a reaffirmation, and a vivification

\item dissolution and debasement
\end{itemize}


From this perspective, the Western nations are moving toward a more gynocratic form. This is not \emph{imposed}, but rather accepted and even chosen, by a race in decline, that is, those who have followed the second possibility. There is still the first possibility.
\end{quotationx}

\paragraph{The Masculine and Feminine Races}


The notion of the ``race of man'' and the ``race of woman'' is not a totally gratuitous extension of the concept. We believe, in fact, that in order to truly take these things into account one should not disregard at all the observations made in Otto Weininger's notable work\footnote{\url{https://www.gornahoor.net/library/SexAndCharacter.pdf}}, especially in two points. The first is in the determination of the type of the pure man or pure woman, as the basis for measuring the ``quantities'' of one or the other that are found in each individual and then to govern oneself in a compliant way. The second is in the daring idea that the existing relations between true man and true woman correspond analogically to those that exist between the Aryan and Semitic races. For Weininger, man is to the woman as the Aryan to the Semite. Weininger did the research on feminine qualities that appear as an exact match to those typical of the Semite and the Jew. One such research is tendentious to a great degree. Weininger, being half-Jewish, was began especially to lose heart and degrade even without wishing it---he did not find the true value of the woman where he should have found it. Nevertheless, his views are still valid; the idea that, from the point of view of a normal and differentiated conception of the sexes, man and woman almost manifest as the expression of two different races, if not pure opposites. It is therefore a serious defect of descriptive and typological racial theory to not take that into account when making the effort to identify and describe the characteristics of each race, while not asking if certain qualities, normal for the masculine type of a given race, continue to be so regarded when applied instead to its feminine type.

In order to remedy this defect, naturally we need to consider the sexes also in their psychical and spiritual aspect. From that point of view, it is certainly absurd to conceive, for example, as normal that the ``Nordic'' woman embodies the same values typical of the pure Nordic man---that is to say, everything that is calm and dominating superiority, solarity, sense of distance, active detachment, connected to a readiness to attack and whatever else we find out. Without directly referring to the Semitic people (unlike Weininger), if we do not want to end up at leveling and therefore mongrelization, it is instead desirable and normal that the same Nordic woman has her own psychical and spiritual qualities, which have a central position in different, non-Nordic races.

Moreover, once the race of the body, or anthropological race, where the female racial differences are well-known and obvious, is omitted, the distinctive characteristics of the women of the various races are much less pronounced on the level of the race of the soul than in the case of the male. In fact, then, at the level of race of the spirit, there is often a true undifferentiation. The real support of the race of the soul and, especially of the spirit, is the man. It is especially in him that the principle of differentiation falls, while the principle of equality is reflected more in the feminine element. Not by accident did the ancient traditions associate, in cosmic analogy, the feminine element with that of matter or unformed power, hylé, dynamis, and the masculine, instead, with the celestial principal of order and individualization. Still less is it an accident that the ancient gynocratic and matriarchal civilizations, as an immediate consequence of the preeminence accorded to the feminine principle in its various, maternal or aphroditic forms, had promiscuity, communism, mother right, general equality as their characteristics. See The Solar Race by Bachofen Strictly speaking, opposite to every man worthy of this name, the true woman that our ancestors considered significantly as ``woman of race'', had to present herself as something dangerous, as a foreign principle that attracts, insinuates herself, and demands an inner action: almost the same type of reaction of which we said speaking of cross-breedings, where foreign blood puts the type to the test and gives rise to two possibilities: either to an awakening, a reaffirmation, and a vivification, or else to dissolution and debasement. In the first case, the man maintains himself at the peak of his purpose and, according to traditional teachings, his qualities permeate and persist intact in his progeny, as the ``dominant'' character. In the second case, in a form more or less disguised, he will undergo at least an internal degeneration of the type. Uncontrollable forces will take the upper hand in the processes of heredity, the protection of the race will become something problematic, until the limit-case is realized, that is, the return, in new forms, to the spirit and promiscuity of the gynocratic civilizations formed by the anti-Nordic races and from the degeneration of the Nordic race.


\flrightit{Posted on 2015-04-29 by Aeneas}

\begin{center}* * *\end{center}

\begin{footnotesize}
\begin{sffamily}
\texttt{Boreas on 2015-04-30 at 10:04 said:}

The Gornahoor wordpress site seems to have been made private; what do I need to do to be able to read the article?

\hfill

\texttt{Cologero on 2015-04-30 at 12:12 said:}

Boreas, we plan to leave it open for certain windows. Eventually, it will probably be made private for those interested. Since only 20\% of readers are clicking through, it really doesn't matter.

\hfill

\texttt{Mark Citadel on 2015-04-30 at 18:11 said:}

Correct me if I'm wrong, is Evola saying here that a `race of the sex' dichotomy is spiritually present during the procreation event and that in Modernity a feminine, degenerative descent is occurring rather than a masculine, heroic descent (that being proper to Tradition.)

If so, what is the solution to this problem, which is obviously not medical or genetic? Will this simply return to the old masculine sexo-racial line of descent once the World of Tradition is restored, or does something active have to be done as it pertains to child-bearing and conception?

\hfill

\texttt{Átila on 2015-05-15 at 11:38 said:}

Read the book ``The Myth of Matriarchal Prehistory'' By CYNTHIA ELLER

\hfill

\texttt{Cologero on 2015-05-15 at 12:35 said:}

Sorry to have to make a public example of you, Atila, but that is not how to write an intelligent comment. If you want to help your fellow human beings, you would have mentioned specific points from the book and then applied them to the post you are commenting on.

I read the amazon page for the book as well as some comments by those who actually read the book. As one reviewer wrote:

``Firstly, it must be understood that Dr. Eller was not out to prove that prehistory was in any way universally patriarchal. Dr. Eller is certainly aware of the many matriarchal, matrilineal and matrifocal prehistorical societies.''

This more or less confirms, or at least is consistent with, Evola's thesis.

\end{sffamily}
\end{footnotesize}

